%******************************************************
%*    LaTex-Vorlage für wissenschaftliche Arbeiten    *
%*    --------------------------------------------    *
%*      Autor: Jan Brennenstuhl (@jbspeakr)           *
%*             www.funkblocka.de                      *
%*                Version: 0.1                        *
%*              Datum: 07.09.2012                     *
%*            Lizenz: CC BY-SA 3.0                    *
%******************************************************
%*  Anmerkung: Unter Verwendung der Paket-            *
%*             Empfehlungen von Georg Verweyen und    *
%*             Teilen der Vorlage von Tino Weinkauf   *
%******************************************************
\documentclass[
  paper=A4, 		% Stellt auf A4-Papier ein
  pagesize, 		% Diese Option reicht die Papiergröße an alle Ausgabeformate weiter
  DIV=calc, 		% Für einen guten Satzspiegel
  headings=small,	% Für etwas kleinere Überschriften
  ngerman,  		% Neue Rechtschreibung (Silbentrennung)
  12pt, 			% Schriftgröße
  listof=totoc, 
  bibliography=totoc, 
  index=totoc, 
  BCOR5mm, 			% Formatregeln für Bindung
]{scrbook}
% ---------------------------------------------------------------
% Unverzichtbare Pakte
\usepackage[T1]{fontenc}	% für Texte mit Umlauten und/oder Akzenten 
\usepackage[utf8]{inputenc}	% Silbentrennung und Eingabe von Umlauten
\usepackage{fixltx2e}  		% Verbessert einige Kernkompetenzen von LaTeX2e
\usepackage[ngerman]{babel} 	% Deutsch als Hauptsprache
\usepackage[ngerman]{translator}
\usepackage{natbib} 
\usepackage{graphicx} 		% Einfügen von Grafiken (.png)
\usepackage[acronym,toc,numberedsection,nonumberlist]{glossaries}
\usepackage[hyphens]{url} 	% URL's schöner formatieren
\usepackage{hyperref}           % Hyperlinks im PDF-Dokument
%\usepackage{pdfpages} 		% Einfügen von PDF-Dateien
% ---------------------------------------------------------------
% Typografisch empfehlenswerte Pakete
\usepackage{
  ellipsis, % Korrigiert den Weißraum um Auslassungspunkte
  ragged2e, % Ermöglicht Flattersatz mit Silbentrennung
  marginnote,% Für bessere Randnotizen mit \marginnote statt % \marginline
}
\usepackage[tracking=true]{microtype}%
\DeclareMicrotypeSet*[tracking]{my}% 
  { font = */*/*/sc/* }% 
\SetTracking{ encoding = *, shape = sc }{ 45 }% alle Passagen in Kapitälchen automatisch leicht gesperrt.
% ---------------------------------------------------------------
% Verschiedene Schriften
\usepackage{%
  lmodern, % A) Latin Modern Fonts sind die Nachfolger von Computer
            % Modern, den LaTeX-Standardfonts
%  hfoldsty % B) Diese Schrift stellt alle Ziffern, außer
            % im Mathemodus, auf Minuskel- oder Mediäval-Ziffern um.
            % Wenn Ihre pdfs unscharf aussehen installieren Sie bitte
            % die cm-super-Fonts (Type1-Fonts).
% charter   % C) Diese Zeile lädt die Charter als Schriftart
}
% ---------------------------------------------------------------
% Glossar-Anpassungen
% Den Punkt am Ende jeder Beschreibung deaktivieren
\renewcommand*{\glspostdescription}{}

% Glossar-Befehle anschalten
\makeglossaries

% Akronyme
%******************************************************
%*    LaTex-Vorlage für wissenschaftliche Arbeiten    *
%*    --------------------------------------------    *
%*      Autor: Jan Brennenstuhl (@jbspeakr)           *
%*             www.funkblocka.de                      *
%*            Lizenz: CC BY-SA 3.0                    *
%******************************************************
\newacronym[\glslongpluralkey={Public-Key-Infrasturkturen}]{PKI}{PKI}{Public-Key-Infrasturktur}
\newacronym{IDM}{IDM}{Identity-Management}
\newacronym{ESB}{ESB}{Enterprise-Service-Bus}
\newacronym{ID}{ID}{Identifikator} 
\newacronym{SigG}{SigG}{Signaturgesetz}
\newacronym{CCC}{CCC}{Chaos Computer Club e. V.}
\newacronym{BSI}{BSI}{Bundesamt für Sicherheit in der Informationstechnik}

% Glossar
%******************************************************
%*    LaTex-Vorlage für wissenschaftliche Arbeiten    *
%*    --------------------------------------------    *
%*      Autor: Jan Brennenstuhl (@jbspeakr)           *
%*             www.funkblocka.de                      *
%*                Version: 0.1                        *
%*              Datum: 07.09.2012                     *
%*            Lizenz: CC BY-SA 3.0                    *
%******************************************************
\newglossaryentry{glos:PKI}{
name=Public-Key-Infrasturktur,
description={Eine Public-Key-Infrasturktur ist eine...}
}
\newglossaryentry{glos:SC}{
name=Smart-Card,
description={EineSmart-Card ist eine...}
}
\newglossaryentry{glos:SigG}{
name=Signaturgesetz,
description={EineSmart-Card ist eine...}
}
% ---------------------------------------------------------------
% Quellen
\bibliography{web,literature}
\usepackage{multibib}
\newcites{literature}{Literaturverzeichnis}
\newcites{web}{Web-Quellenverzeichnis}
% ---------------------------------------------------------------
% Einstellungen
\pdfminorversion=6
%******************************************************
%*    LaTex-Vorlage für wissenschaftliche Arbeiten    *
%*    --------------------------------------------    *
%*      Autor: Jan Brennenstuhl (@jbspeakr)           *
%*             www.funkblocka.de                      *
%*            Lizenz: CC BY-SA 3.0                    *
%******************************************************
%* Anmerkung: der Befehl \hypenation versteht keine   *
%* Sonderzeichen, also weder ä noch "a noch \"a.      *
%* Wörter die derartige Zeichen enthalten müssen      *
%* direkt im Text getrennt werden, z.B. Wör\-ter!     *
%******************************************************
\hyphenation{Ma-nage-ment}
\hyphenation{Ma-nage-ment-agent}
\hyphenation{Ma-nage-ment-agent-en}
\hyphenation{Ma-nage-ment-ar-chi-tek-tur}
\hyphenation{Ma-nage-ment-ar-chi-tek-tu-ren}
\hyphenation{Ma-nage-ment-an-wen-dung}
\hyphenation{Ma-nage-ment-an-wen-dung-en}
\hyphenation{Ma-nage-ment-an-for-der-ung}
\hyphenation{Ma-nage-ment-funk-ti-on}
\hyphenation{Ma-nage-ment-funk-ti-onen}
\hyphenation{Ma-nage-ment-kon-zep-te}
\hyphenation{Ma-nage-ment-res-source}
\hyphenation{Ma-nage-ment-in-for-ma-ti-on}
\hyphenation{Ma-nage-ment-res-sour-cen}
\hyphenation{ma-nage-ment-re-le-vante}
\hyphenation{ma-nage-ment-sy-stem}
\hyphenation{ma-nage-ment-sy-steme}
\hyphenation{Ma-nage-ment-in-stru-men-tie-rung}
\hyphenation{Ma-nage-ment-platt-form}
\hyphenation{Sys-te-men}
\hyphenation{Sys-tem-um-ge-bun-gen}
\hyphenation{Sys-tem-ma-nage-ment}
\hyphenation{DHCP}
\hyphenation{Ma-nage-ment-diszi-plinen}
\hyphenation{System-management-architekturen}
\hyphenation{Verwendungs-nachweise}
\hyphenation{Video-einricht-ungen}
\hyphenation{Res-source}
\hyphenation{Res-sourcen}
\hyphenation{Grund-anwendung}
\hyphenation{Grund-anwendungen}
\hyphenation{Basis-anwendung}
\hyphenation{Core}
\hyphenation{Kom-mu-ni-ka-ti-on}
\hyphenation{De-sign-ent-schei-dung}
\hyphenation{Sprung-ad-res-sen}
\hyphenation{Klas-si-fi-ka-ti-on}
\hyphenation{Schreib-recht}
\hyphenation{Be-nut-zer-zer-ti-fi-kat}
\hyphenation{Bau-stein-ent-wi-ckler}
\hyphenation{ad-mi-ni-stra-ti-ve}
\hyphenation{USB-Spei-cher-me-dien}
\hyphenation{Sys-tem}
\hyphenation{Sys-teme}
\hyphenation{his-to-risch-en}
\hyphenation{PixelPlanet}
\hyphenation{Con-tain-er}
\hyphenation{Con-tain-er-for-mat}
\hyphenation{Con-tain-er-for-mate}
\hyphenation{Bi-blio-thek}
\hyphenation{MyCrypt}
\hyphenation{quell-of-fen}
\hyphenation{quell-of-fene}
\hyphenation{Cam-pus-Ma-nage-ment-Sys-tem-en}
\graphicspath{{./gfx/}}
% ---------------------------------------------------------------
% Befehl für Überschriften mit Untertitel
\newcommand{\Chapter}[2]{\chapter[#1]{#1:\\#2}} 
\newcommand{\Section}[2]{\section[#1]{#1:\\#2}}
% ---------------------------------------------------------------
% Hier startet der Inhalt
\begin{document}
% Meta-Informationen für das PDF
%******************************************************
%*    LaTex-Vorlage für wissenschaftliche Arbeiten    *
%*    --------------------------------------------    *
%* 		  Autor: Jan Brennenstuhl (@jbspeakr)           *
%*						 www.funkblocka.de                      *
%* 		Version: 0.1								                    *
%* 		  Datum: 07.09.2012							                *
%*     Lizenz: CC BY-SA 3.0                           *
%******************************************************
%* Anmerkung: Basierend auf einer Vorlage von         *
%*            Tino Weinkauf.                          *
%******************************************************
\pdfinfo{                               % Zusatzinformationen in PDF-Datei;
                                        % alle Werte sind optional.
    /Author (Dein Name)
    /CreationDate (D:20130310111111)    % Datum der Erstellung
                                        % (D:JJJJMMTThhmmss)
                                        % JJJJ  Jahr
                                        % MM    Monat
                                        % TT    Tag
                                        % hh    Stunden
                                        % mm    Minuten
                                        % ss    Sekunden
                                        %
                                        % Standard: Das aktuelle Datum
                                        %
    /ModDate (D:20130310111111)         % Datum der letzten Modifikation
    /Creator (TeX \& TXC)               % Standard: "TeX"
    /Producer (pdfTeX)                  % Standard: "pdfTeX" + pdftex version
    /Title () 													% Der Titel deiner Arbeit
    /Subject ()                         % Kurzbeschreibung bspw. "Master-Arbeit, Universität Potsdam"
    /Keywords ()                        % Komma-getrennte Liste von Schlagworten
}
% ---------------------------------------------------------------
% Titelblätter und Erklärung
\frontmatter
\pagenumbering{Roman}
%******************************************************
%*    LaTex-Vorlage für wissenschaftliche Arbeiten    *
%*    --------------------------------------------    *
%*      Autor: Jan Brennenstuhl (@jbspeakr)           *
%*             www.funkblocka.de                      *
%*            Lizenz: CC BY-SA 3.0                    *
%******************************************************
% Deckblatt
% -----------------------------------------------------
\thispagestyle{empty}
\begin{center}
		
		% Das Logo der Uni (im gfx-Ordner)
    \includegraphics[width=3cm]{Potsdam_logo}\\
    \vspace{.5cm}
    
    % Der Name deiner Uni
    {\Large \sc Universität Potsdam}\\
    \vspace{1cm}
    
    % Der Name deines Instituts
    {\Huge \sc Institut für Informatik\\[1mm]}
    
    % Der Titel deiner Arbeit (in originaler Sprache)
    \vspace{2cm}
    {\Huge \textbf{Der Titel}}\\
    \vspace*{3mm}
    {\Huge \textbf{deiner Arbeit}}\\
    \vspace*{3mm}
    {\large \textbf{inklusive Untertitel}}\\
    
    % Eine formelle Einleitung
    \vspace{2.0cm}
    {\large Dem Lehrstuhl für}\\
    {\large Komplexe Multimediale Anwendungsarchitekturen}\\
    {\large am Institut für Informatik der Universität Potsdam}\\
    {\large zur Erlangung des akademischen Grades eines}\\
    \vspace{1cm}
    
    % Dein zukünftiger akademischer Grad
    {\Large \sc Master of Science}\\
    \vspace{1cm}
    {\large vorgelegte Master-Thesis von}\\
    %\vspace{0.5cm}
    
		% Deine Anrede
    \Large{Herrn/ Frau Dein Name, akadem. Grad}\\
    
    % Dein aktueller Wohnort
    {\large aus Berlin}\\
    \vspace{5cm} % ggf. je nach Zeilenzahl und Schriftgröße des Titels anpassen
\end{center}

\newpage

% -----------------------------------------------------
% Rückseite Deckblatt
% -----------------------------------------------------
\thispagestyle{empty}
\cleardoublepage

% -----------------------------------------------------
% Erste Seite (Titelblatt)
% -----------------------------------------------------
\thispagestyle{empty}

\begin{center}
    \vspace{.5cm}
    
    % Der Name deiner Uni
    {\Large \sc Universität Potsdam}\\
    \vspace{.5cm}
    
    % Der Name deines Instituts
    {\huge \sc Institut für Informatik\\[1mm]}
    \vspace{1cm}

		% Eine kurze formelle Beschreibung
    {\Large Master-Thesis am Lehrstuhl für}\\
    {\Large Komplexe Multimediale Anwendungsarchitekturen}\\
    \vspace{1.5cm}
    
    % Der Titel deiner Arbeit (in originaler Sprache)
    {\huge \textbf{Der Titel}}\\
    \vspace*{2mm}
    {\huge \textbf{deiner Arbeit}}\\
    \vspace*{2mm}
    {\large \textbf{inklusive Untertitel}}\\
    
    % Der Titel deiner Arbeit (in englischer bzw. deutscher Entsprechung)
    \vspace{1cm}
    {\huge \textbf{The Title}}\\ 
    \vspace*{2mm}
    {\huge \textbf{of your Thesis}}\\
    \vspace*{2mm}
    {\large \textbf{including Subtitle}}\\
    \vspace{1.5cm}
    
    \parbox{1cm}{
      \begin{large}
        \begin{tabbing}
	       	% Beteiligte Personen (Beschreibung jeweils gender-gerecht!)
	        KandidatIn: \hspace{1.5cm} \=Dein Name, B.Sc.\\[2mm]
	    	  Gutachterin: \>Prof. Dr.-Ing. habil. Deine GutachterIn\\
	    		Gutachter:	 \>Dr. Dein Gutachter\\[2mm]
	    		Abgabedatum: \> 1. Januar 1970\\
        \end{tabbing}
      \end{large}
    }\\
    \vspace{1.5mm}
    
		% Das Logo der Uni (im gfx-Ordner)
    \includegraphics[width=2.4cm]{Potsdam_logo}
\end{center}
 
\thispagestyle{empty}
\cleardoublepage
    
% Erklärung (Arbeit selbstständig verfasst)
%******************************************************
%*    LaTex-Vorlage für wissenschaftliche Arbeiten    *
%*    --------------------------------------------    *
%* 		  Autor: Jan Brennenstuhl (@jbspeakr)           *
%*						 www.funkblocka.de                      *
%* 		Version: 0.1								                    *
%* 		  Datum: 07.09.2012							                *
%*     Lizenz: CC BY-SA 3.0                           *
%******************************************************
\begin{large}

\vspace*{2cm}
\noindent
Ich versichere, dass ich die von mir vorgelegte Arbeit selbstständig verfasst habe, dass ich die verwendeten Quellen, Internet-Quellen und Hilfsmittel vollständig angegeben habe und dass ich die Stellen der Arbeit - einschließlich Tabellen, Karten und Abbildungen -, die anderen Werken oder dem Internet im Wortlaut oder dem Sinn nach entnommen sind, in jedem Fall unter Angabe der Quelle als Entlehnung kenntlich gemacht habe.\vspace{2cm}

\noindent
Berlin, den 1. Januar 1970

\vspace{3cm}

\hspace*{6cm}%
\dotfill\\
\hspace*{7cm}%
\textit{(Unterschrift des Kandidaten)}

\end{large}
  
\thispagestyle{empty}
\cleardoublepage

% Zitate
%******************************************************
%*    LaTex-Vorlage für wissenschaftliche Arbeiten    *
%*    --------------------------------------------    *
%* 		  Autor: Jan Brennenstuhl (@jbspeakr)           *
%*						 www.funkblocka.de                      *
%* 		Version: 0.1								                    *
%* 		  Datum: 07.09.2012							                *
%*     Lizenz: CC BY-SA 3.0                           *
%******************************************************
\vspace*{5cm}

\begin{center}
    \textbf{"`Lorem ipsum dolor sit amet, consetetur sadipscing elitr, sed diam nonumy eirmod tempor invidunt ut labore et"'\\}
\end{center}
 Lorem Ipsum (*1961),\\ ehemaliger Chairman of the Board, Präsident und CEO von Lorem Ipsum

\vspace*{1cm}
\noindent 


% Abstract
%******************************************************
%*    LaTex-Vorlage für wissenschaftliche Arbeiten    *
%*    --------------------------------------------    *
%* 		  Autor: Jan Brennenstuhl (@jbspeakr)           *
%*						 www.funkblocka.de                      *
%* 		Version: 0.1								                    *
%* 		  Datum: 07.09.2012							                *
%*     Lizenz: CC BY-SA 3.0                           *
%******************************************************
% Deutsche Zusammenfassung
\begin{center}
    \textbf{Zusammenfassung}
\end{center}
\vspace*{1cm}
% -----------------------------------------------------
\noindent Lorem ipsum dolor sit amet, consetetur sadipscing elitr, sed diam nonumy eirmod tempor invidunt ut labore et dolore magna aliquyam erat, sed diam voluptua. At vero eos et accusam et justo duo dolores et ea rebum. Stet clita kasd gubergren, no sea takimata sanctus est Lorem ipsum dolor sit amet. Lorem ipsum dolor sit amet, consetetur sadipscing elitr, sed diam nonumy eirmod tempor invidunt ut labore et dolore magna aliquyam erat, sed diam voluptua. At vero eos et accusam et justo duo dolores et ea rebum. Stet clita kasd gubergren, no sea takimata sanctus est Lorem ipsum dolor sit amet. Lorem ipsum dolor sit amet, consetetur sadipscing elitr, sed diam nonumy eirmod tempor invidunt ut labore et dolore magna aliquyam erat, sed diam voluptua. At vero eos et accusam et justo duo dolores et ea rebum. Stet clita kasd gubergren, no sea takimata sanctus est Lorem ipsum dolor sit amet.

% -----------------------------------------------------
% Englische Zusammenfassung
\vspace*{1cm}
\begin{center}
    \textbf{Abstract}
\end{center}
\vspace*{1cm}
% -----------------------------------------------------
\noindent Lorem ipsum dolor sit amet, consetetur sadipscing elitr, sed diam nonumy eirmod tempor invidunt ut labore et dolore magna aliquyam erat, sed diam voluptua. At vero eos et accusam et justo duo dolores et ea rebum. Stet clita kasd gubergren, no sea takimata sanctus est Lorem ipsum dolor sit amet. Lorem ipsum dolor sit amet, consetetur sadipscing elitr, sed diam nonumy eirmod tempor invidunt ut labore et dolore magna aliquyam erat, sed diam voluptua. At vero eos et accusam et justo duo dolores et ea rebum. Stet clita kasd gubergren, no sea takimata sanctus est Lorem ipsum dolor sit amet. Lorem ipsum dolor sit amet, consetetur sadipscing elitr, sed diam nonumy eirmod tempor invidunt ut labore et dolore magna aliquyam erat, sed diam voluptua. At vero eos et accusam et justo duo dolores et ea rebum. Stet clita kasd gubergren, no sea takimata sanctus est Lorem ipsum dolor sit amet.

% Inhaltsverzeichnis
\tableofcontents 
% ---------------------------------------------------------------
% Hauptteil - Arabische Seitennummerierung
\mainmatter % die eigentliche Arbeit
\pagenumbering{arabic}

%******************************************************
%*    LaTex-Vorlage für wissenschaftliche Arbeiten    *
%*    --------------------------------------------    *
%*      Autor: Jan Brennenstuhl (@jbspeakr)           *
%*             www.funkblocka.de                      *
%*            Lizenz: CC BY-SA 3.0                    *
%******************************************************
\chapter{Einleitung}
\label{chap:einleitung}
Lorem ipsum dolor \gls{ID} amet, consetetur sadipscing elitr, sed diam nonumy eirmod tempor invidunt ut labore et dolore magna aliquyam erat, sed diam voluptua. At vero eos et accusam et justo duo dolores et ea rebum. Stet clita kasd gubergren, no sea takimata sanctus est Lorem ipsum dolor sit amet. Lorem ipsum dolor sit amet, consetetur sadipscing elitr, sed diam nonumy eirmod tempor invidunt ut labore et dolore magna aliquyam erat, sed diam voluptua. At vero eos et accusam et justo duo dolores et ea rebum. Stet clita kasd gubergren, no sea takimata sanctus est Lorem ipsum dolor sit amet. Lorem ipsum dolor sit amet, consetetur sadipscing elitr, sed diam nonumy eirmod tempor invidunt ut labore et dolore magna aliquyam erat, sed diam voluptua. At vero eos et accusam et justo duo dolores et ea rebum. Stet clita kasd gubergren, no sea takimata sanctus est Lorem ipsum dolor sit amet \citeweb[][]{exmpl:web}.
% -----------------------------------------------------
\section{Motivation}
Lorem ipsum dolor sit amet\footnote{Lorem ipsum dolor sit amet, consetetur \gls{glos:PKI} sadipscing elitr.}, consetetur sadipscing elitr, sed diam nonumy eirmod tempor invidunt ut labore et dolore magna aliquyam erat, sed diam voluptua. At vero eos et accusam et justo duo dolores et ea rebum. Stet clita kasd gubergren, no sea takimata sanctus est Lorem ipsum dolor sit amet. Lorem ipsum dolor sit amet, consetetur sadipscing elitr, sed diam nonumy eirmod tempor invidunt ut labore et dolore magna aliquyam erat, sed diam voluptua. At vero eos et accusam et justo duo dolores et ea rebum. Stet clita kasd gubergren, no sea takimata sanctus est Lorem ipsum dolor sit amet. Lorem ipsum dolor sit amet, consetetur sadipscing elitr, sed diam nonumy eirmod tempor invidunt ut labore et dolore magna aliquyam erat, sed diam voluptua. At vero eos et accusam et justo duo dolores et ea rebum. Stet clita kasd gubergren, no sea takimata sanctus est Lorem ipsum dolor sit amet \citeliterature[vgl.][S. 28ff]{exmpl:buch}.
% -----------------------------------------------------
\section{Anwendungsgebiet}
Lorem ipsum dolor sit amet, consetetur sadipscing elitr, "`sed diam nonumy eirmod tempor invidunt ut labore et dolore magna aliquyam erat, sed diam voluptua"' \citeliterature[][S. 12]{exmpl:buch2}. At vero eos et accusam et justo duo dolores et ea rebum. Stet clita kasd gubergren, no sea takimata sanctus est Lorem ipsum dolor sit amet. Lorem ipsum dolor sit amet, consetetur sadipscing elitr, sed diam nonumy eirmod tempor invidunt ut labore et dolore magna aliquyam erat, sed diam voluptua. At vero eos et accusam et justo duo dolores et ea rebum. Stet clita kasd gubergren, no sea takimata sanctus est Lorem ipsum dolor sit amet. Lorem ipsum dolor sit amet, consetetur sadipscing elitr, sed diam nonumy eirmod tempor invidunt ut labore et dolore magna aliquyam erat, sed diam voluptua. At vero eos et accusam et justo duo dolores et ea rebum. Stet clita kasd gubergren, no sea takimata sanctus est Lorem ipsum dolor sit amet.
% -----------------------------------------------------
\section{Zielstellung}
Lorem ipsum dolor sit amet, \glspl{PKI} consetetur sadipscing elitr, sed diam nonumy eirmod tempor invidunt ut labore et dolore magna aliquyam erat, sed diam voluptua. At vero eos et accusam et justo duo dolores et ea rebum. Stet clita kasd gubergren, no sea takimata sanctus est Lorem ipsum dolor sit amet. Lorem ipsum dolor sit amet, consetetur sadipscing elitr, sed diam nonumy eirmod tempor invidunt ut labore et dolore magna aliquyam erat, sed diam voluptua. At vero eos et accusam et justo duo dolores et ea rebum. Stet clita kasd gubergren, no sea takimata sanctus est Lorem ipsum dolor sit amet. Lorem ipsum dolor sit amet, consetetur sadipscing elitr, sed diam nonumy eirmod tempor invidunt ut labore et dolore magna aliquyam erat, sed diam voluptua. At vero eos et accusam et justo duo dolores et ea rebum. Stet clita kasd gubergren, no sea takimata sanctus est Lorem ipsum dolor sit amet.
% -----------------------------------------------------
\section{Aufbau}
Lorem ipsum dolor sit amet, consetetur sadipscing elitr, sed diam nonumy eirmod tempor invidunt ut labore et dolore magna aliquyam erat, sed diam voluptua. At vero eos et accusam et justo duo dolores et ea rebum. Stet clita kasd gubergren, no sea takimata sanctus est Lorem ipsum dolor sit amet. Lorem ipsum dolor sit amet, consetetur sadipscing elitr, sed diam nonumy eirmod tempor invidunt ut labore et dolore magna aliquyam erat, sed diam voluptua. At vero eos et accusam et justo duo dolores et ea rebum. Stet clita kasd gubergren, no sea takimata sanctus est Lorem ipsum dolor sit amet. Lorem ipsum dolor sit amet, consetetur sadipscing elitr, sed diam nonumy eirmod tempor invidunt ut labore et dolore magna aliquyam erat, sed diam voluptua. At vero eos et accusam et justo duo dolores et ea rebum. Stet clita kasd gubergren, no sea takimata sanctus est Lorem ipsum dolor sit amet.
% -----------------------------------------------------
\endinput 
%******************************************************
%*    LaTex-Vorlage für wissenschaftliche Arbeiten    *
%*    --------------------------------------------    *
%*      Autor: Jan Brennenstuhl (@jbspeakr)           *
%*             www.funkblocka.de                      *
%*            Lizenz: CC BY-SA 3.0                    *
%******************************************************
\chapter{Kapitel 2}
\label{chap:two}
Lorem ipsum dolor sit amet, consetetur sadipscing elitr, sed diam nonumy eirmod tempor invidunt ut labore et dolore magna aliquyam erat, sed diam voluptua. At vero eos et accusam et justo duo dolores et ea rebum. Stet clita kasd gubergren, no sea takimata sanctus est Lorem ipsum dolor sit amet. Lorem ipsum dolor sit amet, consetetur sadipscing elitr, sed diam nonumy eirmod tempor invidunt ut labore et dolore magna aliquyam erat, sed diam voluptua. At vero eos et accusam et justo duo dolores et ea rebum. Stet clita kasd gubergren, no sea takimata sanctus est Lorem ipsum dolor sit amet. Lorem ipsum dolor sit amet, consetetur sadipscing elitr, sed diam nonumy eirmod tempor invidunt ut labore et dolore magna aliquyam erat, sed diam voluptua. At vero eos et accusam et justo duo dolores et ea rebum. Stet clita kasd gubergren, no sea takimata sanctus est Lorem ipsum dolor sit amet.
% -----------------------------------------------------
\endinput 
%******************************************************
%*    LaTex-Vorlage für wissenschaftliche Arbeiten    *
%*    --------------------------------------------    *
%* 		  Autor: Jan Brennenstuhl (@jbspeakr)           *
%*						 www.funkblocka.de                      *
%* 		Version: 0.1								                    *
%* 		  Datum: 07.09.2012							                *
%*     Lizenz: CC BY-SA 3.0                           *
%******************************************************
\chapter{Kapitel 3}
\label{chap:three}
Lorem ipsum dolor sit amet, consetetur sadipscing elitr, sed diam nonumy eirmod tempor invidunt ut labore et dolore magna aliquyam erat, sed diam voluptua. At vero eos et accusam et justo duo dolores et ea rebum. Stet clita kasd gubergren, no sea takimata sanctus est Lorem ipsum dolor sit amet. Lorem ipsum dolor sit amet, consetetur sadipscing elitr, sed diam nonumy eirmod tempor invidunt ut labore et dolore magna aliquyam erat, sed diam voluptua. At vero eos et accusam et justo duo dolores et ea rebum. Stet clita kasd gubergren, no sea takimata sanctus est Lorem ipsum dolor sit amet. Lorem ipsum dolor sit amet, consetetur sadipscing elitr, sed diam nonumy eirmod tempor invidunt ut labore et dolore magna aliquyam erat, sed diam voluptua. At vero eos et accusam et justo duo dolores et ea rebum. Stet clita kasd gubergren, no sea takimata sanctus est Lorem ipsum dolor sit amet.
% -----------------------------------------------------
\endinput 
%******************************************************
%*    LaTex-Vorlage für wissenschaftliche Arbeiten    *
%*    --------------------------------------------    *
%*      Autor: Jan Brennenstuhl (@jbspeakr)           *
%*             www.funkblocka.de                      *
%*            Lizenz: CC BY-SA 3.0                    *
%******************************************************
\chapter{Kapitel 4}
\label{chap:four}
Lorem ipsum dolor sit amet, consetetur sadipscing elitr, sed diam nonumy eirmod tempor invidunt ut labore et dolore magna aliquyam erat, sed diam voluptua. At vero eos et accusam et justo duo dolores et ea rebum. Stet clita kasd gubergren, no sea takimata sanctus est Lorem ipsum dolor sit amet. Lorem ipsum dolor sit amet, consetetur sadipscing elitr, sed diam nonumy eirmod tempor invidunt ut labore et dolore magna aliquyam erat, sed diam voluptua. At vero eos et accusam et justo duo dolores et ea rebum. Stet clita kasd gubergren, no sea takimata sanctus est Lorem ipsum dolor sit amet. Lorem ipsum dolor sit amet, consetetur sadipscing elitr, sed diam nonumy eirmod tempor invidunt ut labore et dolore magna aliquyam erat, sed diam voluptua. At vero eos et accusam et justo duo dolores et ea rebum. Stet clita kasd gubergren, no sea takimata sanctus est Lorem ipsum dolor sit amet.
% -----------------------------------------------------
\endinput 
%******************************************************
%*    LaTex-Vorlage für wissenschaftliche Arbeiten    *
%*    --------------------------------------------    *
%* 		  Autor: Jan Brennenstuhl (@jbspeakr)           *
%*						 www.funkblocka.de                      *
%* 		Version: 0.1								                    *
%* 		  Datum: 07.09.2012							                *
%*     Lizenz: CC BY-SA 3.0                           *
%******************************************************
\chapter{Kapitel 5}
\label{chap:five}
Lorem ipsum dolor sit amet, consetetur sadipscing elitr, sed diam nonumy eirmod tempor invidunt ut labore et dolore magna aliquyam erat, sed diam voluptua. At vero eos et accusam et justo duo dolores et ea rebum. Stet clita kasd gubergren, no sea takimata sanctus est Lorem ipsum dolor sit amet. Lorem ipsum dolor sit amet, consetetur sadipscing elitr, sed diam nonumy eirmod tempor invidunt ut labore et dolore magna aliquyam erat, sed diam voluptua. At vero eos et accusam et justo duo dolores et ea rebum. Stet clita kasd gubergren, no sea takimata sanctus est Lorem ipsum dolor sit amet. Lorem ipsum dolor sit amet, consetetur sadipscing elitr, sed diam nonumy eirmod tempor invidunt ut labore et dolore magna aliquyam erat, sed diam voluptua. At vero eos et accusam et justo duo dolores et ea rebum. Stet clita kasd gubergren, no sea takimata sanctus est Lorem ipsum dolor sit amet.
% -----------------------------------------------------
\endinput 
%******************************************************
%*    LaTex-Vorlage für wissenschaftliche Arbeiten    *
%*    --------------------------------------------    *
%* 		  Autor: Jan Brennenstuhl (@jbspeakr)           *
%*						 www.funkblocka.de                      *
%* 		Version: 0.1								                    *
%* 		  Datum: 07.09.2012							                *
%*     Lizenz: CC BY-SA 3.0                           *
%******************************************************
\chapter{Kapitel 6}
\label{chap:six}
Lorem ipsum dolor sit amet, consetetur sadipscing elitr, sed diam nonumy eirmod tempor invidunt ut labore et dolore magna aliquyam erat, sed diam voluptua. At vero eos et accusam et justo duo dolores et ea rebum. Stet clita kasd gubergren, no sea takimata sanctus est Lorem ipsum dolor sit amet. Lorem ipsum dolor sit amet, consetetur sadipscing elitr, sed diam nonumy eirmod tempor invidunt ut labore et dolore magna aliquyam erat, sed diam voluptua. At vero eos et accusam et justo duo dolores et ea rebum. Stet clita kasd gubergren, no sea takimata sanctus est Lorem ipsum dolor sit amet. Lorem ipsum dolor sit amet, consetetur sadipscing elitr, sed diam nonumy eirmod tempor invidunt ut labore et dolore magna aliquyam erat, sed diam voluptua. At vero eos et accusam et justo duo dolores et ea rebum. Stet clita kasd gubergren, no sea takimata sanctus est Lorem ipsum dolor sit amet.
% -----------------------------------------------------
\endinput 
%******************************************************
%*    LaTex-Vorlage für wissenschaftliche Arbeiten    *
%*    --------------------------------------------    *
%*      Autor: Jan Brennenstuhl (@jbspeakr)           *
%*             www.funkblocka.de                      *
%*            Lizenz: CC BY-SA 3.0                    *
%******************************************************
\chapter{Kapitel 7}
\label{chap:seven}
Lorem ipsum dolor sit amet, consetetur sadipscing elitr, sed diam nonumy eirmod tempor invidunt ut labore et dolore magna aliquyam erat, sed diam voluptua. At vero eos et accusam et justo duo dolores et ea rebum. Stet clita kasd gubergren, no sea takimata sanctus est Lorem ipsum dolor sit amet. Lorem ipsum dolor sit amet, consetetur sadipscing elitr, sed diam nonumy eirmod tempor invidunt ut labore et dolore magna aliquyam erat, sed diam voluptua. At vero eos et accusam et justo duo dolores et ea rebum. Stet clita kasd gubergren, no sea takimata sanctus est Lorem ipsum dolor sit amet. Lorem ipsum dolor sit amet, consetetur sadipscing elitr, sed diam nonumy eirmod tempor invidunt ut labore et dolore magna aliquyam erat, sed diam voluptua. At vero eos et accusam et justo duo dolores et ea rebum. Stet clita kasd gubergren, no sea takimata sanctus est Lorem ipsum dolor sit amet.
% -----------------------------------------------------
\section{Zusammenfassung und erzielte Ergebnisse}
Lorem ipsum dolor sit amet, consetetur sadipscing elitr, sed diam nonumy eirmod tempor invidunt ut labore et dolore magna aliquyam erat, sed diam voluptua. At vero eos et accusam et justo duo dolores et ea rebum. Stet clita kasd gubergren, no sea takimata sanctus est Lorem ipsum dolor sit amet. Lorem ipsum dolor sit amet, consetetur sadipscing elitr, sed diam nonumy eirmod tempor invidunt ut labore et dolore magna aliquyam erat, sed diam voluptua. At vero eos et accusam et justo duo dolores et ea rebum. Stet clita kasd gubergren, no sea takimata sanctus est Lorem ipsum dolor sit amet. Lorem ipsum dolor sit amet, consetetur sadipscing elitr, sed diam nonumy eirmod tempor invidunt ut labore et dolore magna aliquyam erat, sed diam voluptua. At vero eos et accusam et justo duo dolores et ea rebum. Stet clita kasd gubergren, no sea takimata sanctus est Lorem ipsum dolor sit amet.
% -----------------------------------------------------
\section{Ausblick und weiterführende Fragestellungen}
Lorem ipsum dolor sit amet, consetetur sadipscing elitr, sed diam nonumy eirmod tempor invidunt ut labore et dolore magna aliquyam erat, sed diam voluptua. At vero eos et accusam et justo duo dolores et ea rebum. Stet clita kasd gubergren, no sea takimata sanctus est Lorem ipsum dolor sit amet. Lorem ipsum dolor sit amet, consetetur sadipscing elitr, sed diam nonumy eirmod tempor invidunt ut labore et dolore magna aliquyam erat, sed diam voluptua. At vero eos et accusam et justo duo dolores et ea rebum. Stet clita kasd gubergren, no sea takimata sanctus est Lorem ipsum dolor sit amet. Lorem ipsum dolor sit amet, consetetur sadipscing elitr, sed diam nonumy eirmod tempor invidunt ut labore et dolore magna aliquyam erat, sed diam voluptua. At vero eos et accusam et justo duo dolores et ea rebum. Stet clita kasd gubergren, no sea takimata sanctus est Lorem ipsum dolor sit amet.
% -----------------------------------------------------
\endinput 
% ---------------------------------------------------------------
% Schlussteil - Römische Seitennummerierung
\backmatter
\pagenumbering{Roman}
\setcounter{page}{11}
% ---------------------------------------------------------------
% Quellenverzeichnisse
\bibliographystyleliterature{natdin}
\bibliographyliterature{literature}
\bibliographystyleweb{natdin}
\bibliographyweb{web}
% ---------------------------------------------------------------
%Glossar ausgeben
\printglossary[style=altlist,title=Glossar]

%Abkürzungen ausgeben
\deftranslation[to=German]{Acronyms}{Abkürzungsverzeichnis}
\printglossary[type=\acronymtype,style=long]
% ---------------------------------------------------------------
% Anhang
\appendix
%\includepdf[pages=-,fitpaper=true,addtotoc={1,section,0,Diagramm: Vergabe,fig:Vergabedia}]{Pfad-zu-einem-pdf-dokument.pdf}
% ---------------------------------------------------------------
\end{document}
